\section{The Chatnotes arctangent programme: Gaussian torsion and its boundary}
\label{sec:chatnotes-arctangent-torsion}

This section reconstructs the arctangent programme in the directly read raw
chronology at physical lines 9597--11087 of the same 17,437-line Chatnotes
export used in Sections~\ref{sec:chatnotes-coefficient-algebra} and
\ref{sec:chatnotes-weighted-path-algebra}.  The controlling file has 538,005
bytes and SHA-256
\nolinkurl{9a80f53764a9da42c0f8e0b6d19ce61f978980a448a683b9c45f04038716160d}.
The direct-locator audit is
\path{intake/chatnotes/ZN-SYMMETRIES-ARCTANGENT-DIRECT-LOCATOR-AUDIT-05.md}.
The passage is read chronologically: its initial exploratory rhetoric is not
used to infer a theorem, and its later correction controls its earlier
terminology.

At physical lines 9597--9602 the user proposed testing the rational Collatz
prefix pieces against Heegner-field and Machin/arctangent constructions.  The
response then fixed a finite positive exponent word
\[
 w=(a_1,\ldots,a_n),\qquad
 A_k(w)=a_1+\cdots+a_k,
\]
and, at lines 9611--9621, isolated the rational quantities
\begin{equation}
\label{eq:chatnotes-arctangent-prefix-pieces}
 r_k(w)=3^{k-1}2^{-A_k(w)}.
\end{equation}
These are the summands of the prefix coordinate
\(\rho(w)\) in \eqref{eq:weighted-prefix-cocycle}; they are not the terms of
an Euler product.

\subsection{The exact multiplicative torsion test}

At lines 9649--9685 and again at lines 10181--10246, the passage introduced
principal arctangent angles and their Cayley transforms.  With notation chosen
to avoid the existing action \(\theta_{(m,A)}\) and integer \(R(w)\), put
\begin{equation}
\label{eq:chatnotes-arctangent-cayley-data}
 \vartheta_k(w)=\arctan r_k(w),\qquad
 z_k(w)=2^{A_k(w)}+i3^{k-1},\qquad
 u_k(w)=\frac{z_k(w)}{\overline{z_k(w)}}
       =\frac{1+i r_k(w)}{1-i r_k(w)}.
\end{equation}
The arctangent is the principal value, so
\(0<\vartheta_k(w)<\pi/2\) and
\begin{equation}
\label{eq:chatnotes-arctangent-cayley-exponential}
 u_k(w)=e^{2i\vartheta_k(w)}\in\Q(i)^\times.
\end{equation}

\begin{proposition}[repaired torsion formulation]
\label{prop:chatnotes-arctangent-torsion-criterion}
For every \(m=(m_1,\ldots,m_n)\in\Z^n\),
\begin{equation}
\label{eq:chatnotes-arctangent-relation-lattice}
 \sum_{k=1}^n m_k\vartheta_k(w)\in\pi\Q
 \quad\Longleftrightarrow\quad
 \prod_{k=1}^n u_k(w)^{m_k}\in\mu_4,
 \qquad \mu_4=\{1,-1,i,-i\}.
\end{equation}
Consequently the raw question defines the exact lattice
\[
 \operatorname{Rel}_{\pi}(w)=
 \left\{m\in\Z^n:
 \sum_{k=1}^n m_k\vartheta_k(w)\in\pi\Q\right\}.
\]
It is the kernel of the homomorphism
\[
 \Z^n\longrightarrow
 \{x\in\Q(i)^\times:x\overline{x}=1\}/\mu_4,\qquad
 m\longmapsto\left[\prod_{k=1}^n u_k(w)^{m_k}\right].
\]
\end{proposition}

\begin{proof}
Equation~\eqref{eq:chatnotes-arctangent-cayley-exponential} gives
\[
 \prod_k u_k(w)^{m_k}
 =\exp\!\left(2i\sum_km_k\vartheta_k(w)\right).
\]
If the sum is in \(\pi\Q\), the product is a root of unity.  It also lies in
\(\Q(i)^\times\), whose complete root-of-unity group is \(\mu_4\).  Conversely,
if the product is \(i^j\), then
\(2\sum_km_k\vartheta_k(w)\equiv j\pi/2\pmod{2\pi}\); hence the sum lies in
\((\pi/4)\Z\).  Multiplication of the \(u_k\)'s proves the homomorphism and
kernel statements.  Every \(u_k\overline{u_k}=1\), so the displayed codomain
is exact.
\end{proof}

The raw versions wrote \(\mu_\infty\).  The replacement by \(\mu_4\) is not a
change of mathematical object: it records the exact torsion subgroup of the
already printed field \(\Q(i)\).  Calcut's Main Lemma and Corollary~3 give the
same \((\pi/4)\Z\) restriction for rational arctangent sums
\cite{Calcut2009}.

\subsection{The correction internal to the chronology}

The passage initially called the product condition a ``finite logarithmic
unit equation.''  At physical lines 10866--10936 it explicitly corrected that
phrase: a classical unit equation is additive, whereas
\eqref{eq:chatnotes-arctangent-relation-lattice} is a multiplicative torsion
test in \(\Q(i)^\times\).  The same correction says that Baker's theorem may
be applied after a particular logarithmic linear form has been proved
nonzero, but supplies no universal nonvanishing theorem for the arctangent
family.  The critical edition therefore retains the exact criterion and does
not turn its kernel into a theorem of triviality.

This distinction matters already for the legal word \(w=(1,1)\).  Its first
two angles satisfy
\begin{equation}
\label{eq:chatnotes-arctangent-exceptional-seed}
 2\arctan(1/2)+\arctan(3/4)=\frac{\pi}{2}.
\end{equation}
The identity follows from
\((2+i)^2(4+3i)=25i\), and it appears, after division by two, in the
two-term classification of Gasull, Luca, and Varona
\cite[Theorem~1]{GasullLucaVarona2023}.  Thus the strongest possible blanket
independence assertion is false; the exact question is classification of the
relation lattice.

\subsection{Published fixed-instance machinery and the remaining uniform question}

Conway, Radin, and Sadun reconstruct St\o rmer's \(d=1\) theory by unique
factorization in \(\Z[i]\): the arguments of one oriented Gaussian prime above
each rational prime \(p\equiv1\pmod4\), together with \(\pi/4\), give
independent coordinates for rational-tangent angles
\cite[Theorems~1--2 and Section~2]{ConwayRadinSadun1999}.  Hence every fixed
finite instance of \eqref{eq:chatnotes-arctangent-relation-lattice} is
algorithmically decidable by factoring the finitely many Gaussian integers
\(z_k(w)\).  Calcut proves the exact root-of-unity restriction and the
single-angle non-torsion criterion \cite[Main Lemma and
Corollaries~1--3]{Calcut2009}.  Gasull, Luca, and Varona classify their stated
two-term, nonzero-right-hand-side formulas with \(2\)-integer arguments and
derive the relevant zero double-angle family
\cite[Section~2]{GasullLucaVarona2023}.

These sources do not state a uniform rank theorem for the parametric Gaussian
integers
\[
 2^{A_k(w)}+i3^{k-1}\qquad(k=1,\ldots,n).
\]
Exact index and web searches did not locate that formulation.  This negative
search result is not a priority certificate.  It identifies the precise point
at which the source-faithful reconstruction ends and separately authored
research begins.

\begin{nonclaim}
The raw chronology proves neither universal non-torsion nor a Collatz endpoint
statement.  Its initial references to zeta, Heegner fields, modular forms,
monster symmetry, CM arithmetic, transfer operators, and period theory do not
establish a map from any of those structures to the lattice
\(\operatorname{Rel}_{\pi}(w)\).  The fixed-word Gaussian valuation theorem,
the modular non-torsion proof, the exact bounded lattice calculation, and the
all-length rigidity conjecture belong to the separately authored research
companion and are not silently attributed to the raw response.
\end{nonclaim}
